\section{Projekt-Setup und Konfiguration}
Analog zum Frontend erfolgt der Einstieg in die Anwendung über eine definierte Startup-Logik, die in der \texttt{Program.cs} und den Konfigurationsdateien verankert ist.

\subsection{Programm-Initialisierung (Program.cs)}
Die \texttt{Program.cs} ist der Entry Point der ASP.NET Core Anwendung. Hier wird der "Generic Host" gebaut, der Dependency Injection Container befüllt und die HTTP-Request-Pipeline (Middleware) konfiguriert.
\begin{lstlisting}[language={[Sharp]C}, caption={Auszug aus der Middleware-Konfiguration}]
	var builder = WebApplication.CreateBuilder(args);
	
	// 1. Service Registrierung (DI)
	builder.Services.AddApplicationServices(builder.Configuration);
	builder.Services.AddControllers();
	builder.Services.AddEndpointsApiExplorer();
	builder.Services.AddSwaggerGen();
	
	var app = builder.Build();
	
	// 2. HTTP Request Pipeline
	if (app.Environment.IsDevelopment())
	{
		app.UseSwagger();
		app.UseSwaggerUI();
	}
	
	app.UseHttpsRedirection();
	app.UseCors("AllowWebApp");
	app.UseAuthentication();
	app.UseAuthorization();
	app.MapControllers();
	
	app.Run();
\end{lstlisting}

\subsection{Konfiguration (appsettings.json)}
Die Anwendungskonfiguration folgt dem \textit{12-Factor App}-Prinzip: Alle umgebungsabhängigen Parameter werden aus der Datei \texttt{appsettings.json} geladen und können zur Laufzeit durch Umgebungsvariablen überschrieben werden (z.B. in Docker Compose). Das folgende Listing zeigt die vollständige Konfigurationsstruktur (sensible Werte wie Passwörter sind aus Sicherheitsgründen ausgeblendet):

\begin{lstlisting}[language={}   , caption={Konfigurationsstruktur der appsettings.json}]
{
  "ConnectionStrings": {
    "DefaultConnection": "Host=<db-host>;Port=5433;Database=EduShelfDb;..."
  },
  "AIService": {
    "Provider": "Ollama",
    "Endpoint": "http://<ollama-host>:11434",
    "ChatModel": "qwen3:14b",
    "EmbeddingModel": "nomic-embed-text",
    "VisionModel": "moondream:1.8b",
    "ContextLengthLimit": 4096,
    "Prompts": {
      "System": "...",     // Steuert das Grundverhalten des LLMs
      "Intent": "...",    // Anleitung zur JSON-Intent-Erkennung
      "Flashcard": "...", // Anleitung zur Flashcard-Generierung
      "Quiz": "..."       // Anleitung zur Quiz-Generierung
    }
  },
  "Minio": {
    "Endpoint": "<minio-host>:9000",
    "AccessKey": "<key>",
    "SecretKey": "<secret>",
    "BucketName": "edushelf-documents",
    "Secure": false
  },
  "EmailSettings": {
    "SmtpHost": "<smtp-server>",
    "SmtpPort": 587,
    "SmtpUser": "<user>",
    "SmtpPass": "<password>",
    "FromEmail": "<sender>"
  },
  "Seeding": {
    "AdminPassword": "<admin-pw>",
    "StudentPassword": "<student-pw>"
  }
}
\end{lstlisting}

Die wichtigsten Konfigurationsblöcke im Überblick:
\begin{itemize}
    \item \textbf{\texttt{ConnectionStrings}}: Verbindungsstring zur PostgreSQL-Datenbank (Host, Port, Datenbankname, Zugangsdaten).
    \item \textbf{\texttt{AIService}}: Zentrale KI-Konfiguration. \texttt{ChatModel} gibt das Sprachmodell an (z.B. \texttt{qwen3:14b}), \texttt{EmbeddingModel} das Modell für die Vektorerzeugung (\texttt{nomic-embed-text}) und \texttt{VisionModel} das Modell für die Bildanalyse. \texttt{ContextLengthLimit} (Standard: 4096 Token) begrenzt die Kontextfenstergröße. Die \texttt{Prompts}-Untersektion enthält die vollständigen System-Prompts für alle KI-gesteuerten Features als konfigurierbare Strings -- so können diese ohne Codeänderung angepasst werden.
    \item \textbf{\texttt{Minio}}: Verbindungsdaten zum S3-kompatiblen Object Storage (Endpoint, Zugangsdaten, Bucket-Name).
    \item \textbf{\texttt{EmailSettings}}: SMTP-Konfiguration für den Versand von Registrierungs-Bestätigungsmails.
    \item \textbf{\texttt{Seeding}}: Passwörter für vordefinierten Admin- und Student-Benutzer, die beim ersten Start automatisch angelegt werden.
\end{itemize}
